% =============================================================
% Chapitre 2 — État de l'art (~7 pages, resserré — cf. sommaire §7.3)
% Revu le 2026-08-11 (pédagogique + 3 relectures techniques ciblées) :
% corrections appliquées, citations complétées.
% =============================================================
\chapter{État de l'art}
\label{ch:etat-art}

Ce chapitre pose les fondements théoriques mobilisés dans la suite du rapport. Il ne s'agit
pas d'un inventaire exhaustif de la littérature Zero Trust --- un travail de recherche plus
long pourrait lui consacrer un chapitre entier --- mais d'une synthèse ciblée sur les
notions qui conditionnent directement les choix de conception présentés au
chapitre~\ref{ch:conception} : les limites du modèle périmétrique, le modèle Zero Trust, la
sécurisation de la chaîne de livraison logicielle, les approches de détection de menaces, et
la gouvernance de l'infrastructure par le code.

\section{Limites du modèle de sécurité périmétrique}

Deux repères orientent l'ensemble de ce chapitre et seront précisés progressivement : le
modèle \emph{Zero Trust}, qui refuse toute confiance accordée à une entité du seul fait de sa
position réseau, et son corollaire opérationnel, le \emph{principe du moindre privilège} ---
n'accorder à chaque entité que les droits strictement nécessaires à l'action qu'elle exécute.
C'est au regard de ces deux repères que les limites du modèle périmétrique, exposées
ci-dessous, prennent tout leur sens.

Pendant plusieurs décennies, la sécurité des systèmes d'information a reposé sur un modèle
dit \emph{périmétrique} : un ensemble de pare-feux, de zones démilitarisées et de réseaux
privés virtuels délimite une frontière entre un extérieur considéré comme hostile et un
intérieur considéré comme digne de confiance. Une fois cette frontière franchie --- par
exemple via des identifiants volés --- un attaquant se déplace en général librement vers ses
cibles, car les communications internes ne sont pas soumises aux mêmes contrôles que les
communications entrantes.

\textit{Un scénario illustre concrètement cette limite : un poste de travail compromis par
hameçonnage permet à un attaquant de récupérer les identifiants VPN d'un employé ; une fois
connecté, il hérite implicitement de la confiance accordée à toute machine \og interne \fg{},
et peut alors explorer le réseau, identifier des serveurs de bases de données non exposés
publiquement, et y accéder --- non pas parce que ses droits l'y autorisent explicitement, mais
parce qu'aucun contrôle supplémentaire ne s'interpose une fois le périmètre franchi.} Ce n'est
pas un défaut d'implémentation isolé : c'est une conséquence structurelle d'un modèle où la
confiance est accordée une fois, à l'entrée, plutôt que vérifiée à chaque interaction.

Ce modèle repose sur une hypothèse de plus en plus fragile à mesure que les organisations
migrent leurs applications vers des plateformes cloud, où la notion même de \og périmètre
réseau \fg{} perd son sens : les ressources ne sont plus regroupées derrière une frontière
unique, mais distribuées entre des services managés, des fonctions serverless et des bases
de données accessibles par API. C'est exactement le contexte dans lequel s'inscrit MENAL :
une infrastructure entièrement hébergée sur Google Cloud Platform, où aucun \og réseau
interne \fg{} au sens classique n'existe.

\section{Le modèle Zero Trust et la publication NIST SP 800-207}
\label{sec:zero-trust-nist}

Le principe Zero Trust renverse l'hypothèse du modèle périmétrique : la position d'une
entité --- utilisateur, service, machine --- sur le réseau ne lui confère plus aucun crédit
de confiance, et c'est la légitimité de chaque demande, prise individuellement, qui doit
être établie au moment où elle se présente \parencite{kindervag2010}.

La publication de référence sur ce sujet est le \emph{NIST Special Publication 800-207,
Zero Trust Architecture} \parencite{nist800207}, qui formalise l'architecture autour de
trois composants logiques --- un moteur de politique, un administrateur de politique et des
points d'application de politique --- et énonce sept principes (\emph{tenets}) fondateurs.
Trois d'entre eux structurent directement la suite de ce rapport : l'accès à chaque ressource
est accordé sur une base par session, réévaluée à chaque transaction ; l'authentification et
l'autorisation sont strictement appliquées de façon dynamique avant tout accès ; et
l'entreprise surveille en continu l'intégrité et la posture de sécurité de l'ensemble de ses
actifs. Le principe du moindre privilège n'est pas énoncé comme un tenet séparé dans le
document, mais il irrigue plusieurs d'entre eux --- notamment l'octroi d'accès restreint au
strict nécessaire pour accomplir une tâche donnée.

Concrètement, les trois composants logiques de la norme collaborent selon un schéma simple :
un sujet --- utilisateur ou charge de travail --- adresse une demande d'accès à une ressource ;
le \emph{moteur de politique} évalue cette demande au regard de la politique en vigueur et des
signaux disponibles (identité, état de l'équipement, contexte de la requête) et rend une
décision d'octroi ou de refus ; l'\emph{administrateur de politique} traduit cette décision en
instruction technique et établit --- ou refuse d'établir --- le chemin de communication ; le
\emph{point d'application de politique} exécute cette décision en autorisant, surveillant ou
interrompant la connexion. Aucune étape de ce schéma ne repose sur la localisation réseau du
sujet.

En complément de NIST SP~800-207, qui formalise l'architecture cible, le \emph{Zero Trust
Maturity Model} de la CISA \parencite{cisa_zt_maturity_model} propose une grille de progression
par paliers (traditionnel, initial, avancé, optimal) sur cinq piliers --- identité, appareils,
réseau, charges de travail applicatives, données. Ce référentiel offre un repère qualitatif
utile pour situer un projet donné sur une trajectoire de maturité plutôt que sur un état binaire
conforme/non conforme --- une nuance pertinente pour un socle d'infrastructure qui, comme celui
de MENAL, applique les principes Zero Trust progressivement plutôt que d'un seul bloc.

Antérieur de plusieurs années à la formalisation NIST et ayant contribué à façonner le
paradigme Zero Trust au sens large, le projet \emph{BeyondCorp} de Google reste l'expérience
industrielle la plus documentée de ce modèle : il a permis de supprimer le recours au VPN
traditionnel pour l'ensemble des accès internes de l'entreprise, en le remplaçant par une
vérification contextuelle de l'identité et de l'état de l'équipement à chaque requête
\parencite{beyondcorp2014}.

\begin{keybox}[title=Ce que MENAL retient du NIST SP 800-207]
Trois éléments de la norme structurent directement la conception présentée au
chapitre~\ref{ch:conception} --- synthèse de plusieurs tenets et du postulat de compromission
énoncés par le document, non des citations littérales : la vérification explicite à chaque
appel (aucune confiance accordée au seul fait d'être \og dans le réseau \fg{}), le principe du
moindre privilège appliqué par compte de service, et la présomption qu'une partie de
l'infrastructure peut déjà être compromise --- ce qui justifie la journalisation systématique
plutôt qu'une confiance a posteriori.
\end{keybox}

\section{DevSecOps et sécurisation de la chaîne de livraison logicielle}

Le déplacement des contrôles de sécurité vers le début du cycle de développement --- souvent
résumé par l'expression \emph{« shift left »} \parencite{shiftleft_smith2001} --- répond à un
constat simple : plus une vulnérabilité est détectée tard dans le cycle de livraison, plus son
coût de correction est élevé. L'approche DevSecOps consiste à automatiser ces contrôles à
l'intérieur même du pipeline d'intégration et de déploiement, en les transformant en
conditions bloquantes de progression plutôt qu'en vérifications réalisées après coup par une
équipe dédiée. \textit{Concrètement, une porte bloquante refuse par exemple la progression du
pipeline dès qu'une image contient une dépendance présentant une vulnérabilité connue de
sévérité critique, plutôt que de se contenter de journaliser cette vulnérabilité pour une revue
ultérieure.}

Au-delà de ces pratiques ponctuelles de scan, un effort de standardisation plus récent cherche
à formaliser des niveaux de garantie vérifiables pour l'ensemble de la chaîne de construction
logicielle. Le cadre \emph{Supply-chain Levels for Software Artifacts} (SLSA), porté par
l'Open Source Security Foundation, définit une échelle de maturité --- de la simple traçabilité
du processus de build jusqu'à la génération de provenances signées et vérifiables --- qui
structure la manière dont un pipeline peut \emph{prouver}, et non simplement affirmer,
l'intégrité d'un artefact publié \parencite{slsa_framework}.

Plusieurs incidents documentés de la chaîne d'approvisionnement logicielle illustrent que la
chaîne de livraison est elle-même une cible et non un simple canal de déploiement : la
compromission d'un pipeline de build pour y insérer du code malveillant avant distribution
\parencite{cisa_aa20_352a}, la fuite de clés d'accès stockées dans des systèmes de gestion de
version --- un phénomène mesuré à grande échelle sur les dépôts publics
\parencite{gitguardian_sprawl} --- et l'utilisation d'images de base obsolètes héritant de
vulnérabilités déjà connues \parencite{gcp_container_best_practices}. Deux pratiques sont
particulièrement pertinentes pour ce projet : la fédération d'identité pour l'authentification
des pipelines (évitant le stockage de clés statiques de longue durée), et l'analyse
systématique des images et dépendances avant leur publication \parencite{owasp_cicd_top10}.

\section{Détection de menaces : règles statiques et approches sémantiques}

Les systèmes de gestion des informations et des événements de sécurité (SIEM) reposent
historiquement sur des règles de détection statiques, qui comparent des événements
journalisés --- isolément ou en séquence sur une fenêtre glissante --- à des signatures ou à
des motifs connus. Cette approche est efficace pour les menaces déjà répertoriées, mais
atteint ses limites face à des événements qui ne correspondent à aucune règle existante tout
en étant sémantiquement proches d'une technique d'attaque documentée --- par exemple dans le
référentiel MITRE ATT\&CK \parencite{mitre_attack,mitre_attack_design2018}.

Les modèles d'embedding sémantique --- qui transforment un texte en vecteur numérique de
telle sorte que deux textes proches en sens le soient aussi dans l'espace vectoriel
\parencite{devlin_bert2019} --- offrent une approche complémentaire : rapprocher un événement
non couvert par une règle d'une technique d'attaque connue, sur la base d'une similarité de
sens plutôt que d'une correspondance exacte. Un choix de conception important, développé au
chapitre~\ref{ch:conception}, est de retenir un modèle d'embedding \emph{non génératif} : la
sortie est un vecteur déterministe, non un texte produit par un modèle de langage. Ce choix
supprime le vecteur d'attaque par injection de prompt --- aucune sortie textuelle générée à
partir du contenu des journaux ne peut être interprétée comme une instruction par le système
--- sans exclure pour autant une manipulation adversariale du vecteur lui-même, un risque
distinct traité au chapitre~\ref{ch:conception}. Il facilite également, sans la garantir à lui
seul, la reproductibilité d'un résultat d'audit : celle-ci suppose en outre que la version du
modèle et le pipeline de comparaison (seuils de similarité, index de recherche vectorielle)
restent eux-mêmes figés dans le temps.

Dans la pratique, les deux approches ne s'excluent pas. La littérature sur la détection hybride
positionne généralement les règles statiques comme un premier filtre à faible coût de calcul,
et réserve l'évaluation sémantique --- plus coûteuse --- aux seuls événements qui échappent à ce
premier filtre, plutôt qu'à l'ensemble du flux journalisé. Cette architecture en cascade limite
le volume d'événements soumis à l'étape la plus coûteuse, tout en conservant la capacité de
rapprocher un événement inédit d'une technique documentée. Les limites propres à l'approche par
embedding --- faux positifs, faux négatifs, dérive du modèle face à des techniques non vues à
l'entraînement --- sont examinées au chapitre~\ref{ch:conception}, au moment où elles peuvent
être mises en regard des choix de seuils retenus pour MENAL.

\section{Infrastructure as Code et gouvernance}

Décrire l'infrastructure par un code déclaratif et versionné --- plutôt que par des
opérations manuelles dans une console --- permet de la traiter comme un artefact à part
entière : revu avant modification, testé, et reproductible de façon quasi identique
\parencite{iac_morris2020}. Cette pratique répond à deux besoins distincts : la
\emph{reproductibilité} (un environnement peut être recréé dans un état fonctionnellement
équivalent, à partir du seul code) et la \emph{détection de dérive} (tout écart entre l'état
déclaré et l'état réel devient visible au lieu de s'accumuler silencieusement)
\parencite{hashicorp_drift}. \textit{Un exemple concret de dérive : un opérateur modifie
manuellement, via la console web du fournisseur cloud, une règle de pare-feu pour débloquer en
urgence un accès --- une pratique courante lorsqu'un incident presse. Sans détection de dérive,
cette modification reste invisible dans le code source : le prochain déploiement automatisé
peut soit l'écraser silencieusement (perte de la correction), soit, si le code n'est jamais
réappliqué, laisser cette exception non documentée s'accumuler indéfiniment.} Décrire
l'infrastructure de façon déclarative --- en indiquant l'état voulu plutôt que la séquence
d'opérations pour l'atteindre --- rend ce type d'écart détectable par simple comparaison entre
l'état déclaré dans le code et l'état réellement observé sur la plateforme cible.

Cette distinction structure directement le choix, développé au
chapitre~\ref{ch:conception}, de traiter l'ensemble de l'infrastructure de MENAL comme du
code Terraform versionné et revu avant application.

\section{Synthèse et positionnement de MENAL}

Le tableau~\ref{tab:positionnement-etat-art} résume, pour cinq critères directement issus des
sections précédentes, la manière dont chacun des trois modèles couverts par ce chapitre y
répond : le modèle périmétrique traditionnel, une architecture cloud native générique (qui
bénéficie des garanties de la plateforme sans les compléter par une démarche Zero Trust
explicite), et l'architecture Zero Trust retenue pour MENAL.

\begin{table}[htbp]
\centering
\caption{Positionnement de l'architecture MENAL face au modèle périmétrique et au cloud natif générique}
\label{tab:positionnement-etat-art}
\begin{tabular}{@{}P{5.5cm}P{3cm}P{3cm}P{3.3cm}@{}}
\toprule
\textbf{Critère} & \textbf{Architecture périmétrique} & \textbf{Cloud natif générique} & \textbf{Zero Trust MENAL} \\
\midrule
Vérification systématique des accès & Non & Partielle & Oui \\
Moindre privilège, défense en profondeur & Non & Partielle & Oui \\
Segmentation et limitation des mouvements latéraux & Non & Partielle & Oui \\
Sécurité intégrée au cycle de déploiement & Non & Oui & Oui \\
Dashboard SIEM --- détection et investigation d'incidents & Non & Non & Oui \\
\bottomrule
\end{tabular}
\end{table}

La colonne \og cloud natif générique \fg{} illustre un point développé au
\S\ref{sec:zero-trust-nist} : migrer vers le cloud lève une partie des limites du modèle
périmétrique (§2.1) par construction, mais ne suffit pas à lui seul --- sans démarche Zero
Trust explicite, la vérification des accès et la segmentation restent partielles, et la
détection outillée (SIEM) reste absente.

Ces quatre notions --- confiance explicite plutôt qu'implicite, sécurisation de la chaîne de
livraison, détection hybride règles/sémantique, gouvernance de l'infrastructure par le code
--- ne sont pas traitées isolément dans ce projet : elles s'articulent autour d'un principe
transversal détaillé au chapitre~\ref{ch:conception}, celui d'une architecture où chaque
contrôle de sécurité peut être démontré par un test de contournement, et où les décisions
d'architecture sont tracées plutôt qu'implicites.

Aucune de ces quatre notions n'est traitée par MENAL comme un absolu binaire : chacune
correspond à un point d'équilibre choisi, documenté et assumé --- au sens de la grille de
maturité par paliers évoquée au \S\ref{sec:zero-trust-nist} --- plutôt qu'à une conformité
totale à un référentiel donné. C'est ce choix, plutôt qu'un état de l'art exhaustif, que le
chapitre~\ref{ch:conception} traduit en décisions d'architecture concrètes pour MENAL, avant
que le chapitre~\ref{ch:mise-en-oeuvre} n'en détaille la mise en œuvre effective.
